% ======================================================================
% Ontology of Continua -- Core 1.3.3
% Cross-K module: crossk_k3_k4.tex
% Cross-level relations between K3 and K4
% ======================================================================

\subsubsection{Межуровневая структура K3–K4}
\label{sec:crossk-k3-k4}

Переход $K_3 \rightarrow K_4$ обозначает возникновение биологического
континуума из химического. Если $K_3$ поддерживает химическую связность,
реакционные пути и устойчивые кластеры, то $K_4$ вводит:
\begin{itemize}
  \item мембраны и формирование границ,
  \item устойчивые градиенты (ионные, pH, концентрационные, редокс),
  \item осмотические и электрохимические потенциалы,
  \item активные и пассивные транспортные потоки,
  \item метаболические циклы и буферные петли,
  \item ранние носители информации (заряд, структура, шаблонизация),
  \item пороговые условия регулирования и протогомеостаз.
\end{itemize}

Это рождение первых жизненоподобных континуумов: систем, способных поддерживать
неравновесную структуру.

% ----------------------------------------------------------------------
\subsubsection{1. Задействованные уровни и их роли}

\paragraph{K3 (химический континуум).}
$K_3$ предоставляет:
\begin{itemize}
  \item оси состава (виды, заряд, стехиометрия),
  \item потенциалы, определяемые ландшафтом реакций,
  \item реакционные и заряды-переносные потоки $J_3$,
  \item устойчивые кластеры и химические циклы,
  \item отсутствие устойчивых границ и неуправляемых градиентов.
\end{itemize}

\paragraph{K4 (биологический континуум).}
В $K_4$ появляются:
\begin{itemize}
  \item явная граница $\partial\Omega(K_4)$: мембраны и компартменты,
  \item градиенты как оси ($A_{\mathrm{grad}}$),
  \item потенциалы $P_{\mathrm{grad}}, P_{\mathrm{ion}}, P_{\mathrm{redox}}$,
  \item потоки $J_4$ (осмотические, ионопроводящие, активный транспорт, метаболические),
  \item циклы $C_4$ (энергетические, буферные, редокс-циклы, насосы),
  \item пороги $\Theta_4$ для устойчивости мембраны, поддержания градиентов и
        осмотического баланса.
\end{itemize}

Роль перехода K₃→K₄ — стабилизация структуры через контролируемую
неравновесную динамику.

% ----------------------------------------------------------------------
\subsubsection{2. Совместные и наследуемые оси и пороги}

\paragraph{Оси.}
Из $K_3$ в $K_4$ наследуются:
\[
A_{\mathrm{chem}}(K_3) \subset A(K_4).
\]
Новые оси включают:
\begin{itemize}
  \item $A_{\mathrm{grad}}$ — градиенты концентрации, pH, редокс и электричества,
  \item $A_{\mathrm{mem}}$ — кривизна и проницаемость мембран,
  \item $A_{\mathrm{exc}}$ — ось протонейронной активации, связывающая с K₅.
\end{itemize}

\paragraph{Пороги.}
В $K_4$ добавляются семейства порогов:
\begin{itemize}
  \item $\Theta_{\mathrm{mem}}$ — пределы разрыва/проницаемости мембраны,
  \item $\Theta_{\mathrm{grad}}$ — минимальные и максимальные устойчивые градиенты,
  \item $\Theta_{\mathrm{osm}}$ — ограничения осмотического давления ($\Delta\pi$),
  \item $\Theta_{\mathrm{redox}}$ — устойчивость редокс-баланса,
  \item $\Theta_{\mathrm{energy}}$ — пороги метаболической жизнеспособности.
\end{itemize}

Если $\Theta_3$ не выполняется (химическая устойчивость нарушена), биологическая
граница не формируется:
\[
\Theta_3^{\mathrm{death}} \Rightarrow \Omega(K_4)=\varnothing.
\]

% ----------------------------------------------------------------------
\subsubsection{3. Межуровневые потоки, циклы и напряжения}

\paragraph{Потоки.}
Потоки $J_4$ включают все химические потоки $J_3$ и дополнительно:
\begin{itemize}
  \item $J_{\mathrm{osm}}$ — осмотический поток воды,
  \item $J_{\mathrm{ion}}$ — ионный поток через каналы или мембрану,
  \item $J_{\mathrm{pump}}$ — активный транспорт, питаемый $P_{\mathrm{energy}}$,
  \item $J_{\mathrm{redox}}$ — перенос электронов и преобразование энергии,
  \item $J_{\mathrm{metabolic}}$ — сопряжённые реакции и контуры транспорта.
\end{itemize}

Эти потоки задают первый неравновесный динамический режим в иерархии.

\paragraph{Циклы.}
Циклы $C_4$ (энергетические, буферные, насосные, редокс) являются
структурными стабилизаторами:
\[
C_{\mathrm{energy}} : P_{\mathrm{source}} \rightarrow J_{\mathrm{pump}}
\rightarrow \text{восстановление градиента} \rightarrow P_{\mathrm{source}}.
\]

Цикл существует, только если:
\[
\oint_C dA > 0 \quad\text{и}\quad S(C) > 0,
\]
где $S(C)$ — стабильность цикла по определению Core.

\paragraph{Напряжение.}
Межуровневое напряжение $T_{3,4}$ отражает несовместимости между
химической реактивностью и биологической устойчивостью:
\[
T_{3,4} = f(\Theta_{\mathrm{mem}} - P_{\mathrm{chem}},\
           \Theta_{\mathrm{grad}} - P_{\mathrm{grad}},\
           J_{\mathrm{react}} - J_{\mathrm{buffer}}).
\]
Избыточное напряжение разрушает градиенты или мембраны,
сводя $K_4$ обратно к $K_3$.

% ----------------------------------------------------------------------
\subsubsection{4. Условия рождения и исчезновения между уровнями}

\paragraph{Рождение \texorpdfstring{$K_4$}{K_4}.}
Биологический континуум появляется, когда:
\[
T_3 > \Theta_{3,\mathrm{dim}}
\quad \text{и} \quad
A_{\mathrm{grad}}, A_{\mathrm{mem}} \in M_4 \setminus A(K_3),
\]
что соответствует появлению границ и устойчивых градиентов.

Пространство состояний расширяется как:
\[
\Omega(K_4) = \Omega(K_3) \cup \Delta\Omega_{\mathrm{bio}}.
\]

\paragraph{Распространение гибели.}
Если мембрана не может поддерживать градиенты (нарушения
$\Theta_{\mathrm{mem}}$, $\Theta_{\mathrm{grad}}$, $\Theta_{\mathrm{osm}}$), то:
\[
\Omega(K_4) \rightarrow \Omega(K_3),
\]
то есть схлопывание в химический режим.

$K_3$ остаётся устойчивой, если не нарушаются её собственные пороги.

% ----------------------------------------------------------------------
\subsubsection{5. Концептуальные примеры}

\paragraph{Оформление осмотической границы.}
Мембрана возникает как стабильная область $\partial\Omega$, удовлетворяющая:
\[
\Delta\pi = RT(C_{\mathrm{in}} - C_{\mathrm{out}})
\]
в пределах $\Theta_{\mathrm{osm}}$.

\paragraph{Организация на основе градиента.}
Устойчивый градиент $pH$ или ионов определяет новую ось $A_{\mathrm{grad}}$, что
позволяет:
\begin{itemize}
  \item протометаболизм,
  \item векторное преобразование энергии,
  \item селективный транспорт.
\end{itemize}

\paragraph{Прото-эксцитация.}
Электрические градиенты $\Delta V$ позволяют поведение типа протоспайка при:
\[
\Delta V > \Theta_{\mathrm{exc}},
\]
связывая K₄ и K₅ через ранние потоки, похожие на информационные.

Эти примеры иллюстрируют суть перехода K3→K4 в сжатой форме: границы,
градиенты и неравновесные циклы становятся первой устойчивой биологической
организацией.
